\subsection{II)}

{\noindent\large\begin{align*}\displaystyle F(x, y, z) &= \sum(1, 2, 3, 5, 6, 7)\end{align*}}
\subsubsection{Obtener la mínima expresión booleana}
\begin{center}
\begin{tabular}{c|ccc|c}
\rowcolor[HTML]{FFCCC9} 
\cellcolor[HTML]{9AFF99}m & x & y & z & \cellcolor[HTML]{FFFFC7}F(x,y,z) \\ \hline
0                         & 0 & 0 & 0 &                                  \\
1                         & 0 & 0 & 1 & 1                                \\
2                         & 0 & 1 & 0 & 1                                \\
3                         & 0 & 1 & 1 & 1                                \\
4                         & 1 & 0 & 0 &                                  \\
5                         & 1 & 0 & 1 & 1                                \\
6                         & 1 & 1 & 0 & 1                                \\
7                         & 1 & 1 & 1 & 1                               
\end{tabular}
\end{center}
\noindent Los ''1'' de la función pasan, mientras que los ''0'' se los niega. Es una suma de productos de las entradas.
{\large\begin{align*}
    F(x,y,z) &= \overline{x}\cdot\overline{y}\cdot z + \overline{x}\cdot y\cdot\overline{z} + \overline{x}\cdot y\cdot z + x\cdot\overline{y}\cdot z + x\cdot y\cdot\overline{z} + x\cdot y\cdot z
\end{align*}}
\noindent Simplificación por Karnaugh:

\begin{center}
    \begin{karnaugh-map}[2][4][1][$z$][$x\qquad y$]
        \minterms{1,2,3,5,6,7}
        \implicant{1}{5}
        \implicant{2}{7}
    \end{karnaugh-map}
\end{center}
\begin{itemize}[nosep]
    \item Primer bloque ($m_1$, $m_3$, $m_7$ y $m_5$) permanece constante $z$, el resto se sacan.
    \item Segundo bloque ($m_2$, $m_3$, $m_6$ y $m_7$) permanece constante $y$, el resto salen.
    \item La función simplificada queda:
\end{itemize}
{\large\begin{align*}F(x,y,z) &= z + y \end{align*}}

\subsubsection{Simulación}
\imagen[]{10cm}{./imagenes/ejercicio2Falstad.png}
\subsubsection{Descripción en Verilog}
\begin{lstlisting}
module descripcionVerilogEjercicio2(
    input inputX,
    input inputY,
    input inputZ,
    output outputFunction
    );
	 assign outputFunction = inputZ | inputY;

endmodule

\end{lstlisting}

\subsubsection{RTL output}
\imagen[Bloque RTL]{15cm}{./imagenes/rtlEjercicio2Verilog.png}
\imagen[Implementación de las compuertas]{15cm}{./imagenes/rtlCompuertasEjercicio2Verilog.png}

\subsubsection{Simulación temporal}

\noindent Código del módulo Verilog Test Fixture:
\begin{lstlisting}
module verilogTestFixtureEjercicio2;

	// Inputs
	reg inputX;
	reg inputY;
	reg inputZ;

	// Outputs
	wire outputFunction;

	// Instantiate the Unit Under Test (UUT)
	descripcionVerilogEjercicio2 uut (
		.inputX(inputX), 
		.inputY(inputY), 
		.inputZ(inputZ), 
		.outputFunction(outputFunction)
	);

	initial begin
		// Initialize Inputs
		inputX = 0;
		inputY = 0;
		inputZ = 0;

		// Wait 100 ns for global reset to finish
		#100;
      inputX = 0; inputY = 0; inputZ= 1;
		#100;
		inputX = 0; inputY = 1; inputZ= 0;
		#100;
		inputX = 0; inputY = 1; inputZ= 1;
		#100;
		inputX = 1; inputY = 0; inputZ= 0;
		#100;
		inputX = 1; inputY = 0; inputZ= 1;
		#100;
		inputX = 1; inputY = 1; inputZ= 0;
		#100;
		inputX = 1; inputY = 1; inputZ= 1;
		// Add stimulus here

	end
      
endmodule

\end{lstlisting}
\imagen[Simulación temporal]{15cm}{./imagenes/simulacionTemporalEjercicio2.png}
